%% LyX 2.4.4 created this file.  For more info, see https://www.lyx.org/.
%% Do not edit unless you really know what you are doing.
\documentclass[english]{article}
\usepackage[T1]{fontenc}
\usepackage[latin9]{inputenc}
\usepackage{enumitem}
\usepackage{amsmath}
\usepackage{amsthm}
\usepackage{cancel}
\usepackage{geometry}
\geometry{verbose,tmargin=1in,bmargin=1in,lmargin=1in,rmargin=1in}

\makeatletter
%%%%%%%%%%%%%%%%%%%%%%%%%%%%%% Textclass specific LaTeX commands.
\numberwithin{equation}{section}
\numberwithin{figure}{section}
\newlength{\lyxlabelwidth}      % auxiliary length 

%%%%%%%%%%%%%%%%%%%%%%%%%%%%%% User specified LaTeX commands.
\usepackage{fancyhdr}
\usepackage{lastpage}
\pagestyle{fancy}
\usepackage{enumitem}
\usepackage{ifthen}
%sets math fonts to be more readable
\everymath=\expandafter{\the\everymath\displaystyle}
%\DeclareMathSizes{10}{10}{10}{10}
\usepackage{multicol}

\usepackage{tikz}
\usetikzlibrary{plotmarks}
\usepackage{pgfplots}
\pgfplotsset{compat=newest}

\usepackage{calc}
\usepackage[nomessages]{fp}

\usepackage{advdate}    % Advancing/saving dates
\usepackage{datetime}   % Dates formatting
\usepackage{datenumber} % Counters for dates
% Added by lyx2lyx
\setlength{\parskip}{\smallskipamount}
\setlength{\parindent}{0pt}

\makeatother

\usepackage{babel}
\begin{document}
\renewcommand{\author}{Dr.\,James Ashe}

\newcommand{\classNum}{131}

\newcommand{\classSec}{A}

\newcommand{\nameType}{Name:}

\newcommand{\examType}{Quiz 1.1--1.2}

\renewcommand{\date}{\formatdate{28}{8}{2013}}

%%First page's header%%%%%%%%%%%%%%%%%%%%%%
\fancypagestyle{firststyle} %%header settings for first pagr
{
	\fancyhead[L]{MTH \classNum{}(\classSec{}) \author{}\\
	\textbf{\examType{}} \date{}}
	\fancyhead[C]{\textbf{\nameType}}
	\setlength{\headsep}{14pt}
}
%%%%%%%%%%%%%%%%%%%%%%%%%%%%%%%%%%%%%%%%%%

%%Next page's header%%%%%%%%%%%%%%%%%%%%%%%
\setlist{nolistsep}
\lhead{\classNum{}(\classSec{})} 
\chead{\textbf{\nameType}} 
\cfoot{\thepage{}~of~\pageref{LastPage}} 
\setlength{\headsep}{5pt} 
%%%%%%%%%%%%%%%%%%%%%%%%%%%%%%%%%%%%%%%%%%%

%%Various settings; see the preamble for other settings%%%%%
%%\setenumerate[0]{leftmargin=12pt,itemindent=0pt,labelwidth=0pt}
\renewcommand{\labelenumi}{\textbf{\arabic{enumi}.}}
\renewcommand{\labelenumii}{\textbf{\alph{enumii}.}}
\thispagestyle{firststyle}
%%%%%%%%%%%%%%%%%%%%%%%%%%%%%%%%%%%%%%%%%%
\newcommand{\xmax}{10}
\newcommand{\xmin}{-10}
\newcommand{\xstep}{0} %must be xmin+n
\newcommand{\ymax}{10}
\newcommand{\ymin}{-10}
\newcommand{\ystep}{0} %must be ymin+n

\newcommand{\Dspace}{\vspace{1cm}}\textbf{You must show your work to receive credit. }Unless stated
otherwise, each problem is worth 5 points.

\emph{Use the graph to complete the following problems.}\begin{multicols}{2}
\begin{enumerate}[resume]
\item {[}1 pt{]} Graph the point $\left(-3,-5\right)$.\vfill
\item {[}2 pts{]} Use the graph to determine the $x-$int(s) of $f(x)$.\vfill
\item {[}2 pts{]} Use the graph to determine the $y-$int(s) of $f(x)$.
\vfill
\item []\renewcommand{\xmax}{10}
\renewcommand{\xmin}{-10}
\renewcommand{\xstep}{0} %must be xmin+n
\renewcommand{\ymax}{10}
\renewcommand{\ymin}{-10}
\renewcommand{\ystep}{0} %must be ymin+n
%%%%%%%
\begin{tikzpicture} 
\begin{axis}[height=5.5cm, 
	width=5.5cm, 
	axis lines=middle,
	minor x tick num={4},
	minor y tick num={4},
	grid=both,
	%xtick={0,50,...,250},
	%ytick={0,10,20,...,40},
	%axis line style={-latex}, 
	xmin=\xmin,xmax=\xmax, 
	ymin=\ymin,ymax=\ymax,
	%restrict y to domain=-1:10,
	no markers, 
	xlabel={$x$},ylabel={$y$}, 
	%every axis x label/.append style={anchor=north}, 
	%every axis y label/.append style={anchor=east}
	]
	
	%%\addplot[color=red,solid,mark=*] coordinates {(-3,-5)};
	\addplot[domain=\xmin:\xmax,thick,samples=100,draw=red,<->]{-1/2*x+3} node[pos=.65,above] {$f(x)$};
	
\end{axis}
\end{tikzpicture}
\end{enumerate}
\end{multicols}

\emph{Solve the equations.}
\begin{enumerate}[resume]
\item [\textbf{4.}]\setcounter{enumi}{5}$8x-5=11$\vfill
\item $5(x-2)+18=4(x+3)$\vfill
\item $\frac{3x}{2}-x=\frac{x}{10}-\frac{8}{5}$\vfill
\end{enumerate}
\emph{Bonus}
\begin{enumerate}[resume]
\item $\frac{8}{x-4}+\frac{2}{x+5}=\frac{12}{(x+5)(x-4)}$
\end{enumerate}

\end{document}
